\section{Frustration-free regional replacement}

The regional construction of \cite[Definition~5.10]{Arad2016RigorousRG}
preserves the full common kernel. Its gap comparison uses squared norms
and the converse for every unit vector. The rank assertion for arbitrary
interior regions requires correction; an eight-qubit example has rank
seven instead of the stated bound six, as shown in
\cite{gap:alvv16_ff_rank}.

\begin{definition}[Regional parity projections and replacement]
    \label{def:alvv_ff_regional_replacement}
    \lean{TNLean.MPS.AreaLaw.ALVV.restrictedTerm,
        TNLean.MPS.AreaLaw.ALVV.regionalProjection,
        TNLean.MPS.AreaLaw.ALVV.regionalReplacement}
    \leanok
    Let $h_0,\ldots,h_{m-1}$ be orthogonal projections on a real or
    complex Hilbert space. Extend the family by zero when necessary.
    For a finite set $S$ of indices, put $h_i^S=h_i$ for $i\in S$ and
    $h_i^S=0$ otherwise, and define
    \begin{align}
        P_S          & =(I-h_{m-1}^S)\cdots(I-h_0^S),                    &
        \widetilde H & =\sum_{i<m,\ i\notin E\cup O}h_i+(I-P_E)+(I-P_O).
        \label{eq:alvv_ff_regional_replacement}
    \end{align}
    Here the projections commute within $E$ and within $O$.
    In a nearest-neighbor chain, $E,O$ are the even and odd bonds
    contained in the chosen region.
\end{definition}

\begin{lemma}[Projections of the common kernels]
    \label{lem:alvv_ff_common_kernel_layers}
    \lean{TNLean.MPS.AreaLaw.ALVV.complementProd_apply_eq_self_iff,
        TNLean.MPS.AreaLaw.ALVV.isStarProjection_restrictedTerm,
        TNLean.MPS.AreaLaw.ALVV.commute_restrictedTerm,
        TNLean.MPS.AreaLaw.ALVV.isStarProjection_regionalProjection,
        TNLean.MPS.AreaLaw.ALVV.sub_regionalProjection_apply_eq_zero_iff}
    \leanok
    \uses{def:alvv_ff_regional_replacement}
    For a commuting selected family, $P_S$ is an orthogonal projection
    and $(I-P_S)x=0$ if and only if $h_ix=0$ for every $i<m$ in $S$.
    Masking preserves projection and commutation identities.
    The complementary product of any commuting family fixes exactly
    its common kernel.
\end{lemma}
\begin{proof}\leanok
    Products of commuting orthogonal projections are orthogonal
    projections. If every selected term annihilates $x$, each factor
    fixes $x$. Conversely, the detectability estimate for an individual
    term has an empty sum of noncommuting terms, so $h_iP_Sx=0$ for
    each selected $i<m$. Apply this identity to a fixed vector of $P_S$.
\end{proof}

\begin{lemma}[Preservation of the full ground space]
    \label{thm:alvv_ff_regional_kernel}
    \lean{TNLean.MPS.AreaLaw.ALVV.regionalReplacement_apply_eq_zero_iff,
        TNLean.MPS.AreaLaw.ALVV.ker_regionalReplacement,
        TNLean.MPS.AreaLaw.ALVV.regionalReplacement_frustrationFree,
        TNLean.MPS.AreaLaw.ALVV.isPositive_regionalReplacement}
    \leanok
    \uses{def:alvv_ff_regional_replacement}
    The operator $\widetilde H$ in
    (\ref{eq:alvv_ff_regional_replacement}) is positive and satisfies
    \begin{align}
        \ker\widetilde H & =\bigcap_{i<m}\ker h_i
        =\ker\left(\sum_{i<m}h_i\right).
        \label{eq:alvv_ff_regional_kernel}
    \end{align}
    Every vector in this kernel is annihilated by each retained exterior
    term and by each regional term $I-P_E,I-P_O$.
    Thus any nonzero original frustration-free ground space is
    preserved, including its degeneracy.
\end{lemma}
\begin{proof}\leanok
    \uses{lem:alvv_ff_common_kernel_layers}
    A sum of positive operators annihilates a vector exactly when each
    summand does. Apply this to the exterior sum and the two layer
    complements. The layer characterization supplies precisely the
    missing regional kernel conditions. Positivity follows by addition.
\end{proof}

\begin{lemma}[Splitting a commuting layer]
    \label{lem:alvv_ff_layer_split}
    \lean{TNLean.MPS.AreaLaw.ALVV.regionalProjection_mul_complement}
    \leanok
    \uses{def:alvv_ff_regional_replacement}
    If the full family commutes, then $P_SP_{S^c}=P_{\{0,\ldots,m-1\}}$,
    where the complement is taken among $0,\ldots,m-1$.
\end{lemma}
\begin{proof}\leanok
    Induct on the product length. Commute the exterior factor through
    the regional prefix. At each index exactly one of the two masked
    factors is nonidentity, so their product is $I-h_i$.
\end{proof}

\begin{definition}[Collapsed and ordered parity families]
    \label{def:alvv_ff_collapsed_layers}
    \lean{TNLean.MPS.AreaLaw.ALVV.collapsedLayer,
        TNLean.MPS.AreaLaw.ALVV.parityOrderedTerm}
    \leanok
    \uses{def:alvv_ff_regional_replacement}
    Let $e_i,o_i$ be two families, each commuting within itself.
    For a finite regional set $S$, the collapsed even family consists
    of the masked exterior terms $e_i^{S^c}$ for $i<m$, followed by
    $I-P_S^e$ at index $m$, and zero terms afterwards. Define the odd
    family in the same way, using a regional set $T$.
    The ordered original family is
    $a_i=o_i$ for $i<m$ and $a_i=e_{i-m}$ for $i\ge m$.
\end{definition}

\begin{lemma}[Equality of the detectability products]
    \label{lem:alvv_ff_detectability_identity}
    \lean{TNLean.MPS.AreaLaw.ALVV.isStarProjection_collapsedLayer,
        TNLean.MPS.AreaLaw.ALVV.commute_collapsedLayer,
        TNLean.MPS.AreaLaw.ALVV.complementProd_collapsedLayer,
        TNLean.MPS.AreaLaw.ALVV.sum_collapsedLayer,
        TNLean.MPS.AreaLaw.ALVV.isStarProjection_parityOrderedTerm,
        TNLean.MPS.AreaLaw.ALVV.complementProd_parityOrderedTerm,
        TNLean.MPS.AreaLaw.ALVV.sum_parityOrderedTerm}
    \leanok
    \uses{def:alvv_ff_collapsed_layers}
    Each collapsed parity family consists of commuting orthogonal
    projections. Its sum is the retained parity sum plus its regional
    complement. Its complementary product equals that of the original
    parity family. Consequently both ordered two-layer products equal
    $D=P^eP^o$, and the ordered original sum equals
    $H=\sum_{i<m}e_i+\sum_{i<m}o_i$.
\end{lemma}
\begin{proof}\leanok
    \uses{lem:alvv_ff_layer_split}
    The last complementary factor of a collapsed layer is its regional
    projection, and the prefix is the exterior product. Apply the layer
    splitting identity. The sum identity separates the last term.
    For the ordered family, induction first identifies the odd prefix
    and then appends the even factors.
\end{proof}

\begin{lemma}[Converse after regional collapse]
    \label{lem:alvv_ff_collapsed_converse}
    \lean{TNLean.MPS.AreaLaw.ALVV.collapsedLayer_converse}
    \leanok
    \uses{def:alvv_ff_collapsed_layers}
    Let $\widetilde H$ be the sum of the two collapsed parity families.
    Every unit vector $x$ satisfies
    $1-3\langle x,\widetilde Hx\rangle\le\|Dx\|^2$,
    with real parts understood over a complex Hilbert space.
\end{lemma}
\begin{proof}\leanok
    \uses{lem:alvv_ff_detectability_identity}
    Apply the constant-three converse for two commuting projection
    families from \cite{QICLean2026Blueprint} to the collapsed families,
    then substitute their sum and complementary-product identities.
\end{proof}

\begin{lemma}[Orthogonal-complement preservation]
    \label{lem:alvv_ff_orthogonal_preservation}
    \lean{TNLean.MPS.AreaLaw.ALVV.complementProd_mem_orthogonal}
    \leanok
    The complementary product of finitely many orthogonal projections
    preserves the orthogonal complement of their common kernel.
\end{lemma}
\begin{proof}\leanok
    If $h_iz=0$, self-adjointness gives
    $\langle z,h_iy\rangle=\langle h_iz,y\rangle=0$.
    Thus each factor $I-h_i$ preserves the orthogonal complement.
    Induct on the number of factors.
\end{proof}

\begin{lemma}[Gap comparison for two collapsed parity families]
    \label{thm:alvv_ff_collapsed_gap}
    \lean{TNLean.MPS.AreaLaw.ALVV.collapsedLayer_gap}
    \leanok
    \uses{def:alvv_ff_collapsed_layers}
    Suppose every projection among the first $2m$ members of the
    ordered original family fails to commute with at most $g$ others.
    Put $K=\bigcap_{i<2m}\ker a_i$ and suppose
    $\gamma\|y\|^2\le\operatorname{Re}\langle y,Hy\rangle$
    for every $y\in K^\perp$, with $\gamma>0$.
    Every unit vector $x\in K^\perp$ then satisfies
    \begin{align}
        \operatorname{Re}\langle x,\widetilde Hx\rangle
         & \ge\frac{\gamma}{3(\gamma+g^2)}.
        \label{eq:alvv_ff_collapsed_gap}
    \end{align}
    This is the analytic comparison in
    \cite[Theorem~5.14(3)]{Arad2016RigorousRG}; the hypotheses explicitly
    retain the original parity families and their commutation degree.
\end{lemma}
\begin{proof}\leanok
    \uses{lem:alvv_ff_detectability_identity,
        lem:alvv_ff_collapsed_converse,lem:alvv_ff_orthogonal_preservation}
    Detectability gives
    $\operatorname{Re}\langle Dx,HDx\rangle\le g^2(1-\|Dx\|^2)$.
    Orthogonal-complement preservation and the original gap imply
    $(\gamma+g^2)\|Dx\|^2\le g^2$.
    Combine this with the collapsed converse and divide by
    $3(\gamma+g^2)>0$.
\end{proof}

\begin{lemma}[A linear lower bound in the normalized regime]
    \label{cor:alvv_ff_gap_ratio}
    \lean{TNLean.MPS.AreaLaw.ALVV.gap_ratio_le}
    \leanok
    For $0<\gamma\le1$,
    $\gamma/15\le\gamma/[3(\gamma+2^2)]$.
\end{lemma}
\begin{proof}\leanok
    The denominator on the right is at most fifteen; multiply by the
    two positive denominators.
\end{proof}

\begin{definition}[Matrix coordinates and bond endpoints]
    \label{def:alvv_ff_matrix_replacement}
    \lean{TNLean.MPS.AreaLaw.ALVV.bondTermExtension,
        TNLean.MPS.AreaLaw.ALVV.bondTermExtension_supported,
        TNLean.MPS.AreaLaw.ALVV.regionalProjectionMatrix,
        TNLean.MPS.AreaLaw.ALVV.toEuclideanCLM_regionalProjectionMatrix,
        TNLean.MPS.AreaLaw.ALVV.regionalReplacementMatrix,
        TNLean.MPS.AreaLaw.ALVV.toEuclideanCLM_regionalReplacementMatrix,
        TNLean.MPS.AreaLaw.ALVV.evenBonds,
        TNLean.MPS.AreaLaw.ALVV.oddBonds}
    \leanok
    Number the chain sites $0,\ldots,L-1$ and extend the finite list
    of original bond terms by zero. Let $h_i$ be supported on
    $\lbrack i,i+2)$. For a site region $\lbrack a,b+1)$ select exactly the bond indices
    $\lbrack a,b)$, separating their even and odd members into $E,O$.
    The matrix parity products are defined recursively by $P_S(0)=I$ and
    $P_S(m+1)=(I-h_m^S)P_S(m)$. The replacement matrix is the exterior
    sum plus $(I-P_E(m))+(I-P_O(m))$.
    Under the canonical identification of matrices with operators on the
    Euclidean chain space, these are exactly
    (\ref{eq:alvv_ff_regional_replacement}).
    This endpoint convention follows the explicit regional sum in
    \cite[Definition~5.10]{Arad2016RigorousRG}; the final qudit index is
    excluded from its bond-index set, as detailed in
    \cite{gap:alvv16_ff_rank}.
\end{definition}

\begin{lemma}[Support and rank of the parity products]
    \label{lem:alvv_ff_matrix_product_rank}
    \lean{TNLean.MPS.AreaLaw.ALVV.regionalProjectionMatrix_supported,
        TNLean.MPS.AreaLaw.ALVV.rank_regionalProjectionMatrix_le_one,
        TNLean.MPS.AreaLaw.ALVV.rank_regionalProjectionMatrix_le_sq}
    \leanok
    \uses{def:alvv_ff_matrix_replacement}
    Suppose $d\ge1$ and each $h_i$ is supported on $\lbrack i,i+2)$.
    A parity product is supported on any set containing every selected
    bond support. Its operator Schmidt rank at each cut is at most $d^2$.
    If no selected bond crosses that cut, its rank is at most one.
    Projection, positivity and commutation assumptions are unnecessary
    for these rank statements.
\end{lemma}
\begin{proof}\leanok
    Induct on the product length. A factor $I-h_i$ remains supported on
    its bond. Across the cut after the first $k$ sites, only the bond
    $i+1=k$ can cross. Every other factor has a one-term tensor
    decomposition. The crossing factor has a decomposition of length
    at most $d^2$ by expansion in matrix units on its single left site.
    Products of these decompositions give the asserted bounds.
\end{proof}

\begin{theorem}[Corrected rank bound for interval replacement]
    \label{thm:alvv_ff_corrected_rank}
    \lean{TNLean.MPS.AreaLaw.ALVV.rank_regionalReplacementMatrix_le_sq_add_three}
    \leanok
    \uses{def:alvv_ff_matrix_replacement}
    For arbitrary $a,b,m,k$ and $d\ge1$, the regional replacement
    of the first $m$ nearest-neighbor terms satisfies
    \begin{align}
        \operatorname{OSR}_{\lbrack0,k)\mid\lbrack0,k)^c}(\widetilde H)
         & \le d^2+3.
        \label{eq:alvv_ff_corrected_rank}
    \end{align}
    Only the stated local supports are required.
    This corrects the unrestricted rank claim of
    \cite[Theorem~5.14(2)]{Arad2016RigorousRG}.
\end{theorem}
\begin{proof}\leanok
    \uses{lem:alvv_ff_matrix_product_rank}
    At a cut inside the selected bond interval, the exterior sum is
    $H_L\otimes I+I\otimes H_R$. Absorb one identity into each side.
    These two one-sided summands contribute rank at most two.
    Exactly one regional parity product can contain the crossing bond;
    the two products together contribute at most $d^2+1$.
    Thus their difference has rank at most $d^2+3$.

    At a cut outside the region, both regional products lie on the same
    side and can be absorbed, together with $2I$, into that side's
    exterior sum. The other one-sided sum contributes at most one,
    and the at most one exterior crossing term contributes at most
    $d^2$. This gives the sharper bound $d^2+2$ in this case.
\end{proof}

\begin{lemma}[Rank bound for a prefix or suffix]
    \label{thm:alvv_ff_boundary_rank}
    \lean{TNLean.MPS.AreaLaw.ALVV.rank_regionalReplacementMatrix_le_sq_add_two_of_boundary}
    \leanok
    \uses{def:alvv_ff_matrix_replacement}
    Under the same local-support hypotheses, if $a=0$ or $m\le b$,
    then every cut satisfies
    $\operatorname{OSR}(\widetilde H)\le d^2+2$.
\end{lemma}
\begin{proof}\leanok
    \uses{lem:alvv_ff_matrix_product_rank}
    At an interior cut, a prefix has no retained exterior term to the
    left, and a suffix has none to the right. The remaining exterior
    sum, together with $2I$, has rank at most one. Add the two parity
    products, whose combined contribution is at most $d^2+1$.
    At an exterior cut, absorb both regional products into their common
    side exactly as in the proof of
    Theorem~\ref{thm:alvv_ff_corrected_rank}.
\end{proof}

\begin{definition}[The chain operators]
    \label{def:alvv_ff_chain_operators}
    \lean{TNLean.MPS.AreaLaw.ALVV.colorTerm,
        TNLean.MPS.AreaLaw.ALVV.chainTerm,
        TNLean.MPS.AreaLaw.ALVV.chainOperator,
        TNLean.MPS.AreaLaw.ALVV.chainReplacement,
        TNLean.MPS.AreaLaw.ALVV.toEuclideanCLM_chainReplacementMatrix,
        TNLean.MPS.AreaLaw.ALVV.chainOperator_eq_toEuclideanCLM_sum}
    \leanok
    \uses{def:alvv_ff_regional_replacement,def:alvv_ff_matrix_replacement}
    Let $h_i$, $0\le i<m$, be matrices on the chain Hilbert space,
    extended by zero beyond $m$. Write $H=\sum_{i<m}h_i$.
    For $c\in\{0,1\}$ let $h_i^{(c)}=h_i$ when $i\equiv c\pmod2$
    and $h_i^{(c)}=0$ otherwise. The operator replacement
    $\widetilde H$ is the matrix replacement under the canonical
    star-algebra identification of matrices with continuous operators.
    Thus $H$ is exactly the supplied finite bond sum and
    $\widetilde H$ is exactly its regional replacement, retaining all
    exterior bonds.
\end{definition}

\begin{lemma}[Locality supplies commutation degree two]
    \label{lem:alvv_ff_chain_degree}
    \lean{TNLean.MPS.AreaLaw.ALVV.chainTerm_commute_of_not_adjacent,
        TNLean.MPS.AreaLaw.ALVV.commutationDegree_parityOrderedTerm_le_two}
    \leanok
    \uses{def:alvv_ff_chain_operators}
    If $h_i$ is supported on $\lbrack i,i+2)$, then $h_i$ and $h_j$ commute
    unless their bond indices are adjacent. The ordered odd-then-even
    family of length $2m$ consequently has commutation degree at most two.
\end{lemma}
\begin{proof}\leanok
    Distinct nonadjacent bonds have disjoint supports, so their
    matrices commute. Each active member of the ordered family has a
    unique original bond index $p$; its position is $m+p$ for even $p$
    and $p$ for odd $p$. Any noncommuting active member must have
    original index $p-1$ or $p+1$. Hence the filtered set of
    noncommuting positions is contained in a set of cardinality at most
    two. Zero factors commute with every factor.
\end{proof}

\begin{lemma}[Kernel, constraints and gap for the chain]
    \label{lem:alvv_ff_chain_comparison}
    \lean{TNLean.MPS.AreaLaw.ALVV.ker_chainReplacement,
        TNLean.MPS.AreaLaw.ALVV.chainReplacement_frustrationFree,
        TNLean.MPS.AreaLaw.ALVV.chainReplacement_gap,
        TNLean.MPS.AreaLaw.ALVV.chainReplacement_gap_linear}
    \leanok
    \uses{def:alvv_ff_chain_operators}
    Suppose $d\ge1$ and the supplied local terms are orthogonal
    projections supported on their nearest-neighbor bonds.
    Then $\ker\widetilde H=\ker H$, and each vector in this kernel
    is annihilated separately by both regional layer complements and
    by every retained exterior term.
    If $\gamma>0$ and
    $\gamma\|y\|^2\le\operatorname{Re}\langle y,Hy\rangle$ for all
    $y\perp\ker H$, every unit $x\perp\ker H$ satisfies
    \begin{align}
        \operatorname{Re}\langle x,\widetilde Hx\rangle
         & \ge\frac{\gamma}{3(\gamma+4)}.
        \label{eq:alvv_ff_chain_gap}
    \end{align}
    For $0<\gamma\le1$, this lower bound is at least $\gamma/15$.
\end{lemma}
\begin{proof}\leanok
    \uses{thm:alvv_ff_regional_kernel,lem:alvv_ff_chain_degree,
        thm:alvv_ff_collapsed_gap,cor:alvv_ff_gap_ratio}
    Locality gives commutation within each parity, so apply full
    kernel preservation. Positivity of the original terms identifies
    the original kernel with the common kernel of the ordered family.
    Write $P^{(c)}_S$ for the product of
    (\ref{eq:alvv_ff_regional_replacement}) formed from the family
    $h^{(c)}$. The two parity sums equal $H$, and masking each parity by
    the regional bond set gives the two regional products:
    \begin{align}
        \sum_{i<m}h_i^{(0)}+\sum_{i<m}h_i^{(1)} & =H,   &
        P^{(0)}_E                               & =P_E, &
        P^{(1)}_O                               & =P_O.
        \label{eq:alvv_ff_chain_parity_sums}
    \end{align}
    The two exterior parity sums add to the retained exterior sum, so
    the collapsed families reassemble the replacement:
    \begin{align}
        \sum_{i<m}\bigl(h^{(0)}\bigr)_i^{E^c}+(I-P_E)
        +\sum_{i<m}\bigl(h^{(1)}\bigr)_i^{O^c}+(I-P_O)
         & =\widetilde H.
        \label{eq:alvv_ff_chain_collapsed_sum}
    \end{align}
    Apply the collapsed-family gap comparison with the derived degree
    two, and substitute (\ref{eq:alvv_ff_chain_parity_sums}) and
    (\ref{eq:alvv_ff_chain_collapsed_sum}). The scalar estimate gives
    the normalized lower bound.
\end{proof}

\begin{theorem}[Regional replacement of a frustration-free chain]
    \label{thm:alvv_ff_chain_replacement}
    \lean{TNLean.MPS.AreaLaw.ALVV.frustrationFree_regionalReplacement_chain}
    \leanok
    \uses{def:alvv_ff_chain_operators}
    Let $d\ge1$, let the finite local terms be orthogonal
    projections supported on $\lbrack i,i+2)$, and let $\Omega$ be a unit
    vector with $H\Omega=0$.
    Suppose $\gamma>0$ is a quadratic-form gap above the full kernel
    of $H$. For every site region $\lbrack a,b+1)$, replace precisely its
    bonds $\lbrack a,b)$ by the two regional parity complements, retaining
    every exterior term. Then both regional terms are orthogonal
    projections, $\widetilde H$ is positive and frustration free on
    $\Omega$, and $\ker\widetilde H=\ker H$.
    Its operator Schmidt rank at every cut is at most $d^2+3$,
    improved to $d^2+2$ when $a=0$ or $m\le b$.
    Every unit vector orthogonal to the full replacement kernel has
    energy at least $\gamma/[3(\gamma+4)]$.
    In particular, a unique original ground state remains unique.
    This is the chain theorem of
    \cite[Theorem~5.14]{Arad2016RigorousRG} with the rank correction
    documented in \cite{gap:alvv16_ff_rank}; the unrestricted printed
    $d^2+2$ assertion is not claimed.
\end{theorem}
\begin{proof}\leanok
    \uses{lem:alvv_ff_common_kernel_layers,thm:alvv_ff_regional_kernel,
        lem:alvv_ff_chain_comparison,thm:alvv_ff_corrected_rank,
        thm:alvv_ff_boundary_rank}
    Use the derived same-parity commutations to obtain the two regional
    projections and positivity. Full kernel equality places the
    original unit ground vector in the replacement kernel and
    identifies the two orthogonal complements. The chain
    comparison gives every termwise constraint and the explicit gap.
    The matrix-to-operator identity identifies these same operators
    with those in the proved interval and boundary rank bounds.
\end{proof}
